\documentclass[a4paper,10pt]{report}\input{../0.Préambule/Préambule_cours_prof}\begin{document}

\pagestyle{empty}
\newpage
\noindent NOM, Prénom et classe : \dotfill

\noindent \begin{minipage}{11cm}
\section*{Les nombres décimaux}
\addcontentsline{toc}{section}{Les nombres décimaux}

\noindent \textit{Il sera tenu compte dans la notation de la propreté ainsi que de la justification apportée à chacune des réponses.}

\noindent \textit{Le barème est donné à titre indicatif. Il pourra être modifié ultérieurement.}

\noindent \textit{L'usage de la calculatrice est \textbf{interdite}.} \hfill \textit{Durée : 30 minutes}
\end{minipage}
\hspace{6mm}
\begin{minipage}{5cm}
\renewcommand{\arraystretch}{1.5}
\begin{tabular}{|l|r|}
\hline
\textbf{1. Chercher :} & $\hspace{4mm}/\hspace{4mm}$\\\hline
\textbf{2. Modéliser :} & $\hspace{4mm}/\hspace{4mm}$\\\hline
\textbf{3. Représenter :} & $\hspace{4mm}/\hspace{4mm}$\\\hline
\textbf{4. Calculer :} & $\hspace{4mm}/\hspace{4mm}$\\\hline
\textbf{5. Raisonner :} & $\hspace{4mm}/\hspace{4mm}$\\\hline
\textbf{6. Communiquer :}$\hspace{2mm}$ & $\hspace{4mm}/\hspace{4mm}$\\\hline
\end{tabular}
\end{minipage}

\medskip
\paragraph{Exercice 1}: \hfill \textbf{4 points}

\noindent En utilisant des ordres de grandeur, montrer que Thomas se trompe quand il écrit sur sa copie :
\begin{enumerate}[font=\color{exer1}]
\renewcommand{\arraystretch}{0.5}
\begin{tabularx}{\linewidth}{*{2}{>{\arraybackslash}X}}
\item $385,82+52,37+98,651+0,795=5 376,636$&
\item $6 953,6201 – 401,2587= 2941,3614$\\
\end{tabularx}
\end{enumerate}

\medskip\noindent\grandscarreaux{\linewidth}{8cm}

\medskip
\paragraph{Exercice 2}: \hfill \textbf{3 points}

\noindent Calculer chacune des sommes suivantes en posant l'opération :
\begin{enumerate}[font=\color{exer1}]
\renewcommand{\arraystretch}{0.5}
\begin{tabularx}{\linewidth}{*{3}{>{\arraybackslash}X}}
\item $9,42+ 308,7$ &
\item $7,53+56,4+ 38,95$ &
\item $72,6+39,45+1,8$\\
\end{tabularx}
\end{enumerate}

\medskip\noindent\grandscarreaux{\linewidth}{6.4cm}

\newpage
\paragraph{Exercice 3}: \hfill \textbf{5 points}

\noindent Recopier et compléter les égalités suivantes :
\begin{enumerate}[font=\color{exer1}]
\renewcommand{\arraystretch}{0.5}
\begin{tabularx}{\linewidth}{*{3}{>{\arraybackslash}X}}
\item $7,54 \times $ \dots ~\dots ~\dots ~ $ = 75,4$ &
\item $18,5 \times $ \dots ~\dots ~\dots ~ $ = 1,85$ &
\item $0,72 \times $ \dots ~\dots ~\dots ~ $ = 72$ \\
\item $1,052 \times $ \dots ~\dots ~\dots ~ $ = 105,2$ &
\item $0,84 = $\dots ~\dots ~\dots ~ $ \times 0,84$ &\\
\end{tabularx}
\end{enumerate}

\medskip\noindent\grandscarreaux{\linewidth}{5.6cm}

\medskip
\paragraph{Exercice 4}: \hfill \textbf{4 points}

\noindent Pose et effectue les multiplications suivantes.
\begin{enumerate}[font=\color{exer1}]
\renewcommand{\arraystretch}{0.5}
\begin{tabularx}{\linewidth}{*{2}{>{\arraybackslash}X}}
\item $6,2 \times 5,97$& 
\item $7,65 \times 1,32$\\
\end{tabularx}
\end{enumerate}

\medskip\noindent\grandscarreaux{\linewidth}{5.6cm}

\medskip
\paragraph{Exercice 5}: \hfill \textbf{4 points}

\noindent Pour faire une salade de fruits, Capucine a besoin de $0,5$kg de pommes, $750$g de poires, $300$g d'oranges, $0,4$kg de bananes et $1$ citron. Quand elle pèse le tout, elle obtient $2$kg. Combien pèse le citron ?

\medskip\noindent\grandscarreaux{\linewidth}{5.6cm}

\end{document}